% !TeX root = ../main.tex
\section{Уравнения, допускающие понижение порядка}

\task{С. 7:1}{Уравнение, не содержащее искомую функцию}
\sourcepage{16--17}
\[
 xy''+x(y')^2+y'=0.
\]
При $x\ne0$ положим $z=y'$. Тогда $z'+z/x=-z^2$.
Случай $z=0$ даёт все постоянные $y$. При $z\ne0$ подстановка
$u=1/z$ приводит к
\[
 u'-\frac ux=1,\qquad
 \left(\frac ux\right)'=\frac1x,\qquad u=x(C_1+\ln|x|).
\]
Теперь
\[
 y'=\frac1{x(C_1+\ln|x|)},\qquad
 y=\ln|C_1+\ln|x||+C_2.
\]
\answer{$y\equiv C$ либо $y=\ln|C_1+\ln|x||+C_2$,
на интервалах $x\ne0$, $C_1+\ln|x|\ne0$.}
Постоянные решения проходят и через $x=0$. Непостоянные ветви к нулю
конечно не продолжаются: $\ln|C_1+\ln|x||\to+\infty$.
Есть и короткая проверка: при $v=e^y>0$ исходное уравнение
равносильно $(xv')'=0$.

\task{С. 7:28}{Уравнение без явного вхождения независимой переменной}
\sourcepage{17--19}
\[
 2yy''=(y')^2\bigl(3-4y(y')^2\bigr),\qquad y(4)=1,\quad y'(4)=-1.
\]
Вблизи начальной точки $y'$ не равно нулю, поэтому можно положить
$y'=p(y)$. По правилу цепочки $y''=p_y p$; отсюда
\[
 2y p_y=p(3-4yp^2),\qquad
 p_y-\frac3{2y}p=-2p^3.
\]
Замена $z=p^{-2}$, $z_y=-2p^{-3}p_y$ даёт
\[
 z_y+\frac3y z=4,\qquad (y^3z)_y=4y^3,
 \qquad z=y+\frac C{y^3}.
\]
Из $p(1)=-1$ следует $z(1)=1$, значит $C=0$.
По знаку начальной производной выбираем $y'=-1/\sqrt y$:
\[
 \sqrt y\dd y=-\dd x,\qquad
 \frac23y^{3/2}=-x+\frac{14}3.
\]
\answer{$\displaystyle y(x)=\left(7-\frac32x\right)^{2/3}$,
$\displaystyle -\infty<x<\frac{14}3$.}
Скобка положительна на всём указанном интервале. При $x\to14/3-0$
производная стремится к $-\infty$, поэтому продолжить решение как
$C^2$-функцию невозможно. Положительная ветвь $p$ не удовлетворяет
условию $y'(4)=-1$.

\task{С. 7:5}{Два последовательных понижения}
\sourcepage{19--20}
\[
 yy''-(y')^2=y'y^2.
\]
Все $y\equiv c$ являются решениями. На непостоянном участке, где $y\ne0$,
положим $y'=p(y)$ и получим
\[
 y p_y p-p^2=py^2.
\]
При $p\ne0$ делим на $p$:
\[
 p_y-\frac py=y,\qquad
 \left(\frac py\right)_y=1,\qquad p=y^2+ay.
\]
Следовательно, $y'=y^2+ay$. При $y\ne0$ замена $z=1/y$ даёт
\[
 z'+az=-1.
\]
\answer{При $a\ne0$:
\[
 \frac1y=be^{-ax}-\frac1a;
\]
при $a=0$: $\displaystyle \frac1y=b-x$.
Кроме того, $y\equiv0$. Нули знаменателей исключаются.}
При $b=0$ в первом семействе получаются постоянные $y=-a$, поэтому
все ненулевые постоянные уже включены. Ненулевая ветвь не достигает
нуля за конечное время: это видно из явных формул. Следовательно,
дополнительных склеек с $y=0$ нет.
\editnote{В рукописи в промежуточной строке после $z=1/y$ ошибочно
указано $z'-az=1$. Верное уравнение --- $z'+az=-1$;
конечная формула в рукописи соответствует именно ему.}

\task{С. 7:46}{Однородность по искомой функции}
\sourcepage{20--22}
\[
 yy''+yy'\tan x+2(y')^2=0,\qquad \cos x\ne0.
\]
Сначала воспроизведём метод рукописи. При $y\ne0$ положим $z=y'/y$.
Тогда $y''=y(z'+z^2)$, откуда
\[
 z'+z\tan x=-3z^2.
\]
Решение $z=0$ соответствует постоянным $y$. При $z\ne0$ замена
$w=1/z$ даёт $w'-w\tan x=3$. Умножая на $\cos x$, получаем
\[
 (w\cos x)'=3\cos x,\qquad w=\frac{3\sin x+C_1}{\cos x}.
\]
Таким образом, $y'/y=\cos x/(3\sin x+C_1)$.
Чтобы сразу включить постоянные решения, удобно переписать ответ как
\answer{$\displaystyle y^3=A\sin x+B$, то есть
$\displaystyle y=\sqrt[3]{A\sin x+B}$, на интервалах $\cos x\ne0$
и гладкости этой функции. При $A=0$ получаются все постоянные.}
Проверка без деления на $y$: для $v=y^3$ имеем
\[
 v''+v'\tan x=3y\bigl(yy''+2(y')^2+yy'\tan x\bigr)=0.
\]
Если $A\ne0$ и $\cos x\ne0$, нуль $A\sin x+B$ простой, так что
кубический корень там не имеет конечной производной. Такой нуль
исключается. В рукописной форме $C_2\sqrt[3]{\sin x+C_1}$ ненулевые
постоянные решения теряются.

\task{С. 7:65, а}{Масштабная замена}
\sourcepage{22--23}
\[
 x^2y''+2x^2yy'+2xy^2-2y=0,\qquad y(1)=2,\quad y'(1)=0.
\]
Уравнение сохраняет вид при масштабе $x\mapsto\lambda x$,
$y\mapsto\lambda^{-1}y$. Поэтому для $x>0$ вводим
$t=\ln x$, $y=z(t)e^{-t}$. Точки над $z$ ниже обозначают производные по $t$:
\[
 y'=e^{-2t}(\dot z-z),\qquad
 y''=e^{-3t}(\ddot z-3\dot z+2z).
\]
Подставляя и сокращая $e^{-t}$, находим
\[
 \ddot z-3\dot z+2z\dot z=0,
 \qquad \frac{\dd}{\dd t}(\dot z-3z+z^2)=0.
\]
Начальные данные переходят в $z(0)=2$, $\dot z(0)=2$, поэтому
постоянная интегрирования равна $2-6+4=0$:
\[
 \dot z=3z-z^2.
\]
При $u=1/z$ получаем $\dot u+3u=1$,
$u=1/3+Ke^{-3t}$. Условие $u(0)=1/2$ даёт $K=1/6$.
\answer{$\displaystyle y(x)=\frac{6x^2}{1+2x^3}$, $x>0$.}
Это равносильно ответу рукописи $y/(xy-3)=-2x^2$.
Рациональная формула также даёт гладкое продолжение через $x=0$
на $(-\sqrt[3]{1/2},+\infty)$ и непосредственно удовлетворяет исходному
ненормализованному уравнению. В точке $x=0$ стандартная теорема
единственности для нормальной формы неприменима; саму возможность
продолжения здесь проверяет явная формула.
\editnote{Знак перед $3\dot z$ в одной промежуточной строке рукописи
исправлен: должно быть $-3\dot z$, а не $+3\dot z$.}

\task{Ф. 505}{Понижение порядка и начальные данные}
\sourcepage{24--25}
\[
 y''\cos y+(y')^2\sin y=y',\qquad
 y(-1)=\frac\pi6,\quad y'(-1)=2.
\]
В начальной точке $y'\ne0$ и $\cos y\ne0$.
Полагая $y'=p(y)$, после деления на $p$ получаем
\[
 p_y\cos y+p\sin y=1,\qquad
 \left(\frac p{\cos y}\right)_y=\frac1{\cos^2 y}.
\]
Отсюда $p=\sin y+C\cos y$. Начальная производная определяет
\[
 2=\frac12+C\frac{\sqrt3}2,\qquad C=\sqrt3.
\]
Значит, $y'=2\cos(y-\pi/6)$. Обозначим $v=y-\pi/6$ и интегрируем:
\[
 \ln\tan\left(\frac\pi4+\frac v2\right)=2(x+1),
 \qquad v(-1)=0.
\]
На выбранной ветви $-\pi/2<v<\pi/2$, так что тангенс положителен.
\answer{$\displaystyle y(x)=2\arctan e^{2(x+1)}-\frac\pi3$,
$x\in\R$.}
В частности, $y'=2\cos v>0$ и $y''=-2\sin v\,y'$.
Подстановка $y'=\sin y+\sqrt3\cos y$ даёт
\[
 y''\cos y+(y')^2\sin y
 =y'\bigl[(\cos y-\sqrt3\sin y)\cos y
 +(\sin y+\sqrt3\cos y)\sin y\bigr]=y'.
\]
Эта проверка верна и при $\cos y=0$, хотя промежуточное деление там
было недопустимо. Ответ рукописи с логарифмом равносилен нашему
только после выбора ветви, соответствующей начальным данным.

\task{Д. 2}{Однородность по функции и уравнение Бернулли}
\sourcepage{26--27}
\[
 xy^2y''+x^2(y')^3-xy(y')^2-15y^2y'=0,
 \qquad y(1)=1,\quad y'(1)=4.
\]
Вблизи $x=1$ вводим $z=y'/y$; тогда $y''=y(z'+z^2)$.
После деления на $xy^3$ члены $z^2$ сокращаются:
\[
 z'-\frac{15}x z=-xz^3.
\]
Здесь $z(1)=4$. Подстановка $u=z^{-2}$ приводит к
\[
 u'+\frac{30}x u=2x,\qquad
 (x^{30}u)'=2x^{31},\qquad
 u=\frac{x^2}{16}+Cx^{-30}.
\]
Из $u(1)=1/16$ следует $C=0$. На $x>0$ начальное значение $z(1)=4$
выбирает $z=4/x$, а не $-4/x$.
\[
 \frac{y'}y=\frac4x,\qquad y=Kx^4,\qquad K=1.
\]
\answer{$y=x^4$ на $x>0$; та же формула является решением на всей $\R$.}
Прямая подстановка даёт коэффициент
$12+64-16-60=0$ при $x^{11}$. Глобальная единственность продолжения
через нуль не следует из локальной теоремы: например, можно взять
$y=Kx^4$ при $x<0$ и $y=x^4$ при $x\ge0$ с любым $K$.
Это $C^2$-решение исходного уравнения с теми же данными при $x=1$.
